\section{Introduction}\label{sec:intro}

A renewable stock must be kept above a floor by a manager who knows the structure of the dynamics but not the signs of its parameters. If there are $m$ uncertain signs, the parameter space is the vertex set $\{\pm1\}^m$ of a cube, and the manager's information at any time is a subset of those vertices---the \emph{information class} of cells that the observations so far cannot distinguish. Within a class the manager must act \emph{blindly}: one action for every cell of the class, because the true cell is unknown. The companion obstruction calculus \cite{Abaee2026calc} shows that this is where viability under incomplete observation fails even when every cell is individually controllable: there may be no common action, or no common action \emph{sequence}, that is safe for the whole class. The present article studies the finite, exactly computable instance of that phenomenon on the parameter cube, and answers three questions completely for small cubes: which sets of cells can be kept alive by one blind sequence; from which initial stock; and what the observation of the stock buys.

\paragraph{The model.}
A cell is a sign vector $\theta\in\{\pm1\}^m$; an action is a sign vector $u\in\{\pm1\}^m$ (the manager's guess of the signs, which sets the instrument) or a hold. Per period the stock $z$ moves by $d(u,\theta)=-\tfr12+\tfr15\langle u,\theta\rangle$ and by $-\tfr12$ under the hold; the floor is $z\ge1$. Measured in tenths of stock above the floor the drift is the integer $D(u,\theta)=2\langle u,\theta\rangle-5=2m-5-4\dH(u,\theta)$, a function of the Hamming distance between guess and truth alone, and the initial condition is the integer slack $L=10(z_0-1)$. A set $S$ of cells is \emph{blind-survivable at level $L$} if one action sequence keeps every cell of $S$ at or above the floor for ever when all of them start at slack $L$. The question is simple to state and surprisingly rich: the family of survivable sets is hereditary and monotone in $L$, so it is described by its antichain of maximal members at each level, and these antichains change shape with $m$.

\paragraph{What the earlier version established and what was missing.}
The previous editions of this work (through v19) assembled exact results for $m=4$ and $m=5$: the pair law, the absence of survivable triangles at $m=4$, the two bands of the four-parameter problem, the radius-one balls and the squares of the five-cube, point-based evaluations on $17^4$ action words, and a Lean layer. Referee reports on v19 identified structural problems that this edition addresses: the ``census'' baseline was the power set of the cells (an artefact, not a comparison); the pairwise bound was called sharp without saying from which level; the deadline law was proved only for one schedule; the observation model was named but never computed; the point-based proposition was tautological as stated; the five-cube classification relied on a periodic search that cannot certify the absence of larger survivable sets; the six-cube was left open with a count of $10{,}752$ maximal cliques; and the paper was organised as a technical report. This edition is a reconstruction around one theorem, with the remaining material rebuilt on top of it and all numbers regenerated by a single exact module.

\paragraph{Contributions.}
\begin{enumerate}[leftmargin=2em,itemsep=1pt]
\item \emph{A linear criterion} (Theorem~\ref{thm:criterion}). $S$ is blind-survivable from some finite slack if and only if
\[V(S)=\max_{y\in[-1,1]^m}\ \min_{\theta\in S}\ \langle\theta,y\rangle\ \ge\ \tfr52,\qquad\text{equivalently}\qquad \min_{q\in\Delta(S)}\Big\|\sum_{\theta\in S}q_\theta\theta\Big\|_1\ \ge\ \tfr52 .\]
A maximiser $y$ is a rational mixed action whose product decomposition gives an explicit periodic word; a minimiser $q$ is a pooled-drift certificate of non-survivability, of support at most $m+1$. A non-survivable set loses a cell within $\lfloor L/(5-2V(S))\rfloor$ steps (Corollary~\ref{cor:horizon}), which recovers the far-pair horizon $10(z_0-1)$ of the Lean layer and the deadline law, and is attained.
\item \emph{Structural corollaries.} Pairs are survivable iff $\dH\le m-\tfr52$; the sets survivable at the floor are exactly the subsets of Hamming balls of radius $\lfloor(2m-5)/4\rfloor$; every survivable set has diameter at most $\lfloor m-\tfr52\rfloor$ and hence, by Kleitman's theorem, at most $2,6,12,29,58$ cells for $m=4,\dots,8$, a bound attained by survivable sets for every $m$; and the instruments are incomparable: the tetrahedron of the five-cube has $\sum_i|y_i|=2$ for the uniform certificate although each of its cells is individually kept by an action of full strength (Section~\ref{sec:criterion}).
\item \emph{Thresholds with two-sided certificates} (Section~\ref{sec:thresholds}). The least sufficient slack $\Lstar(S)$ is the least $L$ with $(L,\dots,L)$ in the winning region of a safety game on $\N^S$; the region is an up-set, computed as an antichain of minimal winning vectors. A saturated version of the game gives a sound upper bound, exact forward reachability with dominance pruning gives the lower bound, and the two met for all $5{,}373$ representative sets of this paper. The first-step bound $\min_u\max_{\theta\in S}(-D(u,\theta))^+$ is the threshold for most, but not all, maximal sets.
\item \emph{Complete classifications for $m=4,5,6$} (Section~\ref{sec:classification}). The six-cube is new: the $10{,}752$ maximal cliques of its distance-three graph fall into seven orbits of the hyperoctahedral group, three of which are survivable; the maximal survivable sets are $192$ Kleitman-extremal twelve-cell cylinders ($\Lstar=1$), $1{,}280$ ten-cell sets on the boundary $V=\tfr52$ ($\Lstar=3$) and $160$ three-dimensional subcubes ($\Lstar=5$). The level census counts $7{,}072$ survivable sets at the floor, $631{,}440$ from slack $1$, $774{,}800$ from slack $3$ and $774{,}960$ from slack $5$, against $853{,}680$ pairwise-admissible candidates; the $2{,}400$ minimal obstructions inside the candidate family all have $V=2$.
\item \emph{Observation computed, not named} (Section~\ref{sec:observation}). The blind value of a belief is the largest mass of a survivable set, which gives the probe ladders for $m=4,5,6$; the \emph{feedback game}, in which the observed stock reveals the drift and splits the class, is solved exactly ($\Fstar=23$ for the four-cube, $18$ for the five-cube, with the shell and subcube thresholds); the first splitting action costs the antipode $2m+5$ tenths; and total ignorance for $T$ periods costs exactly $5T$ tenths whatever the sequence.
\item \emph{Point-based evaluation placed correctly} (Section~\ref{sec:pbvi}). The alpha-set of $k$-step outcome masks is exact for $k$-step survival and its value at a belief is a maximum over at most $857$ ($m=4$) or $15{,}929$ ($m=5$) masks; but it overestimates indefinite survival until $k>L$, because far pairs survive exactly $L$ steps. Probed beliefs (a fixed seed, adversarial priors) show the gap.
\item \emph{Why exact arithmetic} (Section~\ref{sec:float}). $143{,}360$ survivable sets of the six-cube have $V=\tfr52$ exactly, and their optimal words return to the floor every period. A float32 simulation declares a false collapse at step $3$; float64 gives $1{,}120$ false collapses among the $1{,}280$ maximal boundary sets under one coding of the drift and none under another. Runtimes of the whole exact pipeline are reported.
\item \emph{Verification} (Section~\ref{sec:verification}). One exact module regenerates every number ($\NCHECKS$ checks); the Lean layer is $8$ modules and $80$ theorems with no \texttt{sorry}, and its scope is stated precisely.
\end{enumerate}

\paragraph{Relation to the companion articles.}
The obstruction calculus \cite{Abaee2026calc} defines the common-action obstruction and the information classes; the certificate article \cite{Abaee2026cert} builds continuous-to-finite bridges with exact witnesses; the minimax article \cite{Abaee2026minimax} gives the dual (measure) certificate for one-step obstructions on general action sets; the probabilistic article \cite{Abaee2026prob} introduces belief-state safety values and the survivable-set antichains. The present article is the finite-cube laboratory of that programme: the one-step measure certificate of \cite{Abaee2026minimax} becomes, over sequences, the pooled-drift certificate $q$ of Theorem~\ref{thm:criterion}, and the antichain formula of \cite{Abaee2026prob} is the device of Section~\ref{sec:observation}.

\paragraph{Organisation.}
Section~\ref{sec:model} fixes the model and the information classes. Section~\ref{sec:criterion} proves the linear criterion and its corollaries. Section~\ref{sec:thresholds} defines and certifies thresholds. Section~\ref{sec:classification} states the classifications and the census. Section~\ref{sec:observation} treats observation. Section~\ref{sec:pbvi} places point-based evaluation. Section~\ref{sec:float} compares exact and floating-point computation and reports runtimes. Section~\ref{sec:verification} is the verification record and Lean scope. Section~\ref{sec:related} discusses related work and open problems. Proofs omitted from the main text, all records and tables, and the complete check list are in the supplementary material (sections S1--S9).
